\par\noindent\begin{minipage}{\linewidth}
% !TeX root = ../教材题库_全部条目_转写工作稿.tex
% Source: XJ-B1-C01-EX-P022-Y0891
% 第1章 / 1.2 常用逻辑用语 / 教材 p.22 / PDF 第29页
% 已对照原页转写题干及公式；待教师复核；未制作答案
% 原题及相邻原解截图（核对用）：../题目原图/XJ-B1-C01-EX-P022-Y0891-01.webp
% 本文件只在转写完成后作为题干使用；OCR 条目尚不是可靠题干。
\begin{bankitem}{1.2 例题 例题9}{第1章 / 1.2 常用逻辑用语 / 教材 p.22 / PDF 第29页}{XJ-B1-C01-EX-P022-Y0891}{已对照原页转写题干及公式；待教师复核；未制作答案}
写出下列含量词命题的否定：
\subq{1}$p$：所有二次函数的图象都是轴对称图形；
\subq{2}$q$：$\exists x\in\mathbb R,\ -2x^2+4x-3>0$；
\subq{3}$r$：三角形的垂心都在其外部；
\subq{4}$s$：有一个小于 210 的正整数至少有 4 个质因数。
\end{bankitem}
\end{minipage}\par\medskip
